\scanpage{85}
les $a_j$ qui sont dans $\widetilde{\mathcal T}_s$ forment une famille génératrice des flèches de $\widetilde{\mathcal T}_s$ --- ce qui implique que tous les sommets de $\widetilde{\mathcal T}$ sont parmi \deleted{\unsure{des}} les origines et extrémités des $a_j$.
\deleted{De l'hypothèse \unsure{ici}}
Supposons que $\forall\alpha$, les \deleted{\unsure{origines}\ \unsure{sommets et les}\inserted{\unsure{paires}} des flèches $\widetilde u'_{i'_{\alpha,\nu}}$ $(0\leq\nu\leq n_\alpha)$ de $\widetilde{\mathcal T}$},
\[\{\widetilde t'_{i'_n},\widetilde s'_{i'_{n-1}}\}\ \cdots\ \{\widetilde t'_{i'_2},\widetilde s'_{i'_1}\}\quad\text{dans (9)},\]
\deleted{qui sont}\inserted{couples}, qui (par composabilité de la suite \deleted{\unsure{dans}}\inserted{des $u'_{i'}$} figurant dans $\unsure{R}$), sont contenus chacun dans un même groupoïde fibre de $\widetilde{\mathcal T}\text{---}\mathcal T$, soient contenus dans une même composante connexe de la catégorie fibre correspondante.

L'hyp. sur les \unsure{$a_j$} implique donc qu'il existe un mot \deleted{\unsure{g}} composable dans $\widetilde{\mathcal T}$ en les $\unsure{a_j^\varepsilon}$, qui joigne $\widetilde t_{i_n}\unsure{\,:}\widetilde s_{i_{n-1}}$, et
\[(11)\]
\diagram{\begin{tikzcd}[ampersand replacement=\&, column sep=large]
\deleted{\unsure{\widetilde t_{i'_n}}}\widetilde s'_{i'_n}\arrow[r,"u'_{i'_n}"]\&\widetilde t'_{i'_n}\arrow[r,squiggly,"M_{\alpha,i_n}((\unsure{a_j^\varepsilon}))"]\&\widetilde s'_{i'_{n-1}}\arrow[r,"u'_{i'_{n-1}}"]\&\widetilde t'_{i'_{n-1}}\arrow[r,squiggly,"M_{\alpha,i_{n-1}}((a_j))"]\&{}\\
\&\cdots\&\widetilde s'_{i'_2}\arrow[r]\&\widetilde t'_{i'_2}\arrow[r,squiggly,"M_{\alpha,2}((a_j))"]\&\widetilde s'_{i'_1}\arrow[r,"u'_{i'_1}"]\&\widetilde t'_{i'_1}
\end{tikzcd}}
\mspage{84}
